\begin{savequote}
\begin{verbatim}
 _________________________________________
( "Why make it simple when it is possible )
( to make it complicated"                 )
(                                         )
( Husse March 2 2007                      )
 -----------------------------------------
       o   ,__,
        o  (oo)____
           (__)    )\
              ||--|| *

\end{verbatim}
\end{savequote}
\chapter
[\CO{[MP1]} Medizinprodukte 1]
{\CC{[MP1]}\\Medizinprodukte 1}
\label{chp:GER}
\label{chp:MP1}

\ChapterIntro
{%
Das folgende Kapitel soll einen Überblick über die
wichtigsten Funktionen der entsprechenden Geräte bieten. Es
ersetzt keine vom Hersteller herausgegebene
Gebrauchsanweisung. Diese ist vom Betreiber dem Personal in
geeigneter Form zugänglich zu machen und muss von den
Anwendern konsultiert werden.
}
{%
\begin{description}
  \item[Maintainer:] Roman Koch
  \item[Autoren:] Michael Neuhold
  \item[Reviewer:] --
  \item[Version:] {\DocVersion} ({\DocVersionInfo})
  \item[SHA1:] {\DocGitCommit}
\end{description}

{\TxTeilkapitelAASS}
}


\clearpage % TODO temp clearpage
\Layout{0}{\TxQuerverweise}
{~

\noindent\includegraphics[angle=270,width=\linewidth]{./images/leitsystem/Leitschema-PAM_MP_.pdf}
}{\begin{itemize}
  \item \EE{Beschreibung der Medizinprodukte}: \CO{[MP1, MP2]} 
  (\ref{chp:MP1}, \prettyref{chp:MP2}).
  \item \EE{Untersuchungen}: \CO{[BAU]} (\prettyref{chp:BAU}).
  \item \EE{Sanitätshilfemaßnahmen}:  \CO{[SAN]} (\prettyref{chp:SAN}).
%   \item[Anamnese und Patientengespräch:]  Kapitel [AUP] (\ref{chp:AUP}).
\end{itemize}}{}{}{}{}



\section{Medizinproduktegesetz - MPG}
\label{topic:mpg-info}

\Layout{61}{}{%
\noindent Das Medizinproduktegesetz regelt die Funktionstüchtigkeit,
Sicherheit, Inbetriebnahme, Instandhaltung, den Betrieb und
die Anwendung, Überwachung, Sterilisation, Desinfektion und
Reinigung von Medizinprodukten. Als Medizinprodukte gelten
alle Instrumente, Apparate, Vorrichtungen, Stoffe oder
andere Gegenstände (inkl. Software), die vom Hersteller zur
Anwendung am Menschen zwecks Erkennen, Überwachen und Behandlung von
Krankheiten, Verletzungen oder Behinderungen bestimmt sind\ZFootnote{Genauer:
Zwecks
\begin{inparaenum}
  \item Erkennung, Verhütung, Überwachung, Behandlung oder Linderung von Krankheiten,
  \item Erkennung, Überwachung, Behandlung, Linderung oder Kompensierung von Verletzungen oder Behinderungen,
  \item Untersuchung, Veränderung oder zum Ersatz des anatomischen Aufbaus oder physiologischer Vorgänge oder
  \item Empfängnisregelung,
\end{inparaenum}
wobei \ldots{} Medikamente fallen nicht darunter.}.

Gegenstände und Stoffe, die dazu bestimmt sind, zusammen mit einem
Medizinprodukt verwendet zu werden ("`Zubehör"'), gelten selber auch als
Medizinprodukt. Dies trifft auch auf \E{Desinfektionsmittel} zu, mit welchen
Medizinprodukte gereinigt und wiederaufbereitet werden!


Es gibt tausende verschiedene Produkte, i.\,d.\,R. haben sie
alle eine \E{CE-Kennzeichnung}. Diese zeigt an, dass das
Produkt gewissen Minimalanforderungen entspricht. 
}
{}
{./images/geraete/Ce-logo.jpg}
{Das CE-Kennzeichen.}
{---}




\Layout{57}{}
{}{}
{
\begin{tabulary}{\linewidth}{lLL}
%\addlinespace
\TabHeading{Klasse} &
\TabHeading{Beschreibung} &
\TabHeading{Beispiele}
\\\addlinespace\midrule
\TabSub{Klasse I}
&
Kein oder geringes Anwendungsrisiko
&
z.B: Verbandsmaterial, Gehhilfe, Rollstuhl, Tragsessel, Sehhilfe, Stethoskop, 
Trage, Beatmungsbeutel, Fieberthermometer, Druckminderer 
\\\addlinespace
\TabSub{Klasse IIa}
&
Anwendungsrisiko ohne entsprechender Ausbildung bzw. Unterweisung
&
z.B: aktive Diagnosegeräte, Absaugeinheit und -katheter, EKG-Kabel, 
Zuckermessstreifen, Trachealtuben, Dauerkatheter 
\\\addlinespace
\TabSub{Klasse IIb}
&
Erhöhtes Anwendungsrisiko ohne entsprechender Ausbildung bzw. Unterweisung
&
z.B: Beatmungsgerät, Defibrillator, Pulsoxymeter, Babyinkubator, Röntgengeräte
\\\addlinespace
\TabSub{Klasse III}
&
Besonders hohes Anwendungsrisiko ohne entsprechende Ausbildung bzw. Unterweisung
&
z.B: Perfusor, Insulinpumpe, Herzkatheter oder -klappen, Medizinprodukte mit Arzneimittelkomponenten
\\\addlinespace\bottomrule
\end{tabulary}
}
{MPG-Risikoklassen.}{}

\Layout{0}{Pflichten gemäß MPG}{
Anwender sind verpflichtet, die
Medizinprodukte richtig, zweckentsprechend und sicher einzusetzen bzw.
anzuwenden, so wie es ihnen bei ihrer Einweisung gelehrt wurde und es der
Hersteller vorgesehen hat. Darüber hinaus ergeben sich noch weitere spezielle
Verpflichtungen:

\subparagraph{Funktionskontrolle} 
Der Anwender muss sich vor jeder Anwendung eines
Medizinproduktes von der Funktionstüchtigkeit,
Betriebssicherheit und dem ordnungsgemäßen Zustand überzeugen
("`Gerätecheck"').

\subparagraph{Einweisung} 
Medizinprodukte dürfen nur von solchen Personen angewendet
werden, die auf Grund ihrer Ausbildung,  und
erforderlichenfalls einer produktspezifischen \EE{Einweisung}
das Medizinprodukt sachgerecht anwenden können und auch auf
besondere anwendungs- und medizinproduktespezifische Gefahren
hingewiesen worden sind. Dabei sind die Gebrauchsanweisungen
sowie sonstige beigefügten Informationen der Produkte zu
beachten und den Anwendern zugänglich zu machen.

\subparagraph{Meldepflicht} 
Fachpersonal muss Informationen über Medizinprodukte im
Hinblick auf \E{Zwischenfälle}\footnote{Insbesonders jede
Fehlfunktion, aber auch z.\,B. jeden Mangel in Bezug auf die
Kennzeichnung oder die Gebrauchsanweisung, die zum Tod oder
zu einer schwerwiegenden Schädigung eines Patienten,
Anwenders oder eines Dritten führen können oder geführt
haben}, oder schwerwiegende \E{Nebenwirkungen}, oder
schwerwiegende \E{Qualitätsmängel}, unverzüglich dem
Bundesamt für Sicherheit im Gesundheitswesen
melden.\ZFootnote{\S\,70 MPG}
}
{
\begin{itemize}
  \item Sachgerechte Anwendung
  \item Einweisung
  \item Funktionskontrolle
  \item Meldepflicht bei schweren Mängeln
\end{itemize}
}
{}
{}
{}







%%%%%%%%%%%%%%%%%%%%%%%%%%%%%%%%%%%%%%%%%%%%%%%%%%%%%%%%%%%
%%%%%%%%%%%%%%%%%%%%%%%%%%%%%%%%%%%%%%%%%%%%%%%%%%%%%%%%%%%
%%%%%%%%%%%%%%%%%%%%%%%%%%%%%%%%%%%%%%%%%%%%%%%%%%%%%%%%%%%



\subsection{Steriles Material}

\Layout{22}{Entpacken von sterilem Material}
{
Die Verpackung von sterilem Material muss \E{fachkundig}
geöffnet werden, damit es beim Öffnen nicht verunreinigt
wird. Gängiges Verpackungsmaterial lässt sich
\EE{aufschälen}, d.\,h. die Oberseite lässt sich von der
Unterseite ablösen. Dabei darf das Material nicht
kontaminiert werden.

Manche vorgefertigten Sets (z.\,B. für diverse chirurgische
Eingriffe) werden nicht in foliierten Verpackungen
ausgeliefert, sondern sind in mehrere Lagen von Spezialpapier
eingewickelt. Das Spezialpapier darf nur an der Außenseite
angefasst werden und wird auf einer Fläche ausgebreitet,
sodass das darin eingewickelte Material steril bleibt.
}
{
\begin{itemize}
  \item Aufschälen
  \item Material muss steril bleiben!
\end{itemize}
}
{./images/AASdev20090228-img104.jpg}
{Aufschälen einer Verpackung.}
{}



% \begin{figure}[H]
% \includegraphics[width=6.656cm,height=4.993cm]{images/AASdev20090228-img103.jpg}
% \includegraphics[width=6.663cm,height=4.984cm]{images/AASdev20090228-img104.jpg}
% \Captionof{figure}{Öffnen von sterilen Material unter verschiedenen
% Bedingungen: Immer wird die Verpackung "`aufgeschält"'.
% \AuthNotesPicLegalNotOk{}{}}
% \end{figure}

% In diesem Abschnitt werden weitere Maßnahmen und Geräte für die Diagnostik und
% Patientenüberwachung (Monitoring) besprochen. Manche Maßnahmen werden in anderen
% Abschnitten besprochen (z.\,B. Blutdruckmessung
% (\prettyref{topic:blutdruckmessung}), Messung der Herzfrequenz (Pulsmessung,
% \prettyref{topic:pulsmessung}), \ldots). 

\clearpage % TODO temp clearpage

\section{Spezielle Medizinprodukte}

Für Medizinprodukte sind die Medizinprodukteeinweisungen
sowie die Gebrauchsanleitungen bzw. Herstellerhinweise zu
beachten. Allgemeine, nicht produktspezifische Informationen
finden sich in folgenden Abschnitten:

\begin{itemize}
  \item Absauggeräte
  \item Beatmungshilfsmittel
  \item Berieselungshilfsmittel
  \item Sauerstoff
  \begin{itemize}
    \item 
  \end{itemize}
  \item Geräte zum Retten von Patienten
  und Schienungsmaterial
  \begin{itemize}
    \item Aluminiumkernschiene
    \item HWS-Schiene
    \item Rettungskorsett
    \item Schaufeltrage
    \item Spineboard
    \item Vakuummatratze
    \item Vakuukschiene
  \end{itemize}
  \item Transporthilfsmittel
  \begin{itemize}
    \item (Fahr-)Trage
    \item Tragering
    \item Tragetuch
    \item Tragsessel
  \end{itemize}
  \item 
  \item Diagnostik und Monitoring
  \begin{itemize}
    \item 
    \item Blutdruckmessung
    \item Blutzuckermessung
    \item Pulsoxymetrie
    \item Stethoskop
    \item Temperaturmessung
  \end{itemize}
\end{itemize}
\section{Diagnostik und Monitoring}






\subsection{Manuelle Blutdruckmessgeräte}

\subsection{Pulsoxymetrie}


\subsubsection{Nellcor{\protect\texttrademark} NPB-40,
Nellcor{\protect\texttrademark} NPB-40 Oximax,
Nellcor{\protect\texttrademark} N-65 Oximax}


Diese Geräte sind \E{nicht} defibrillationssicher. Sie dürfen während der
Defibrillation am Patienten bleiben, doch während und kurz nach dem Schock können
die Messwerte ungenau sein und das Gerät darf während des Schocks nicht berührt
werden.

\TODO{GABS}{2010-06-14}{POI- Pulsoxy Nellcor: Symbol
"`Batterie"' und "`durchgestrichener Lautsprecher"' einfügen.}{}
Eine leer werdende Batterie wird durch ein Batterie-Symbol gekennzeichnet. Die
leeren Batterien fachgerecht entsorgen und durch 4 neue Batterien der Type "`AA"'
unter Beachtung der Polarität (siehe Symbole im Batteriefach auf der Rückseite
des Gerätes) ersetzen. Die Lautstärke des Pulstones kann durch die Tasten \Sign{↑} und \Sign{↓}
eingestellt werden.

Bei Auftreten eines Alarmes kann dieser für max. 120 Sekunden 
\A{durch Drücken der Taste []}
stummgeschalten werden.




\subsubsection{Masimo{\protect\texttrademark} Rad5v}

Dieses Gerät darf während der Defibrillation verwendet werden, es kann jedoch während des
Schocks und 20 Sekunden danach zu verfälschten Messwerten kommen.
Die Batteriezustandsanzeige ist unterhalb der Ein-/Aus-Taste in Form von 4 grünen LEDs
angebracht. Wenn nur mehr eine grüne LED der Batteriezustandsanzeige angezeigt wird, sind
die leeren Batterien fachgerecht zu entsorgen und durch 4 neue Batterien der Type „AA“
unter Beachtung der Polarität (siehe Symbole im Batteriefach auf der Rückseite des Gerätes)
ersetzen.

\TODO{GABS}{2010-06-14}{Pulsoxy Masimo Rad5v: Symbol "`Taste Pulston"'
einfügen.}{}
\opt{STATUS-DRAFT}{\A{Draft:}
Die Lautstärke des Pulstones kann durch Drücken der Taste [] in 4 Stufen eingestellt werden
(Aus, Laut, Mittel, Leise).
}
Dieses Gerät alarmiert nicht bei niedriger Sauerstoffsättigung, lediglich bei fehlendem
Messsignal und eignet sich somit nur bedingt für Überwachungsaufgaben.


\subsection{Blutzuckermessung}

Die Blutzuckermessung wird unter
\prettyref{topic:blutzucker-messung} besprochen.

\subsubsection{AccuCheck{\protect\texttrademark} go}
% TODO : Neufassung
\TODO{GABS}{yyyy-mm-dd}{AccuCheck{\protect\texttrademark}
go: Neufassung.}{}

{\TakeNotesLarge}

\subsection{Fieberthermometer}

% TODO INHALT: Neufassung GER Fieberthermeter
\TODO{GABS}{yyyy-mm-dd}{Fieberthermometer: Neufassung.}{}


\Layout{2}{Ohrthermometer}
{~

\TakeNotesLarge

\subparagraph{Hygienische Maßnahmen}~

\TakeNotesMedium
}{}{}{}{}

\Layout{2}{Standardthermometer}
{~

\TakeNotesLarge

\subparagraph{Hygienische Maßnahmen}~

\TakeNotesMedium
}{}{}{}{}



\section{Transporthilfsmittel}

\Layout{2}{{\TxQuerverweise}}{%
\begin{inparaitem}
  \item Tragen mit dem Tragering: \prettyref{topic:tragering};
  \item Umgang mit dem Rollstuhl: \prettyref{topic:rollstuhl};
  \item Tragen mit dem Tragetuch: \prettyref{topic:tragetuch-taetigkeit}
\end{inparaitem}
}{}{}{}{}

\subsection{Tragsessel}

\Layout{22}{Umgang mit dem Tragsessel}
{\label{topic:tragsessel}%
Der Patient ist immer anzuschnallen. Beim Transport über
Stiegen hat der Patient immer talwärts zu blicken.

{\TakeNotesLarge}
}{

}
{./images/noauth/GER-tragsessel.jpg}
{\AuthNotesPicLegalNotOk{GER-tragsessel.jpg}{}Tragsessel}
{}

\Layout{2}{Zutragestuhl}
{\label{topic:zutragestuhl-mp}%
{\TakeNotesLarge}
}{

}
{}
{}
{}


\subsection{Krankentragen}
\label{topic:fahrtrage}

\Layout{2}{Krankentrage mit Fahrgestell}
{
\index{Trage}
\index{Krankentrage}
\index{Fahrtrage}
\index{Stollenwerk{\texttrademark}}
\index{Ferno!Fahrtrage}
Der Patient ist immer anzuschnallen (Schultern, Becken und Beine!). Beim
Transport über Stiegen hat der Patient immer talwärts zu blicken. 

Wenn ein Patient auf der Trage liegt darf immer nur
\E{\EE{ein} Ende} gleichzeitig gehoben oder gesenkt werden,
da sonst die Gefahr des Kippens besteht.

\opt{ORGASBFLO}{Eingesetzte Geräte:

\begin{itemize}
    \item Hersteller: Fa. Stollenwerk
    \item Hersteller: Fa. Ferno
\end{itemize}
}
}{

}
{}
{}
{}

\subsubsection{Krankentrage mit Fahrgestell der Fa. Stollenwerk}
% TODO : Neufassung
\TODO{GABS}{yyyy-mm-dd}{Krankentrage Stollenwerk: Fehlender
Text.}{Hier muss neuer Text geschrieben werden!}


\Layout{2}{\TxBeschreibung}
{%
\begin{compactdesc}
    \item[Hersteller:] Stollenwerk
    \item[Produktfamilie:] Krankentrage, Fahrgestell
    \item[Modell/Typ:] Trage 3002, 3003, 3006; Fahrgestell 2870
    \item[Modellbesonderheiten:] \begin{inparadesc}
    \item[Trage 3002:] Kopfhochstellung. 
    \item[Trage 3003:] Kopf- und Fußhochstellung. 
    \item[Trage 3006:] Kopf- und Fußhochstellung, Bauchdecken-Entlastung. 
    \item[Alle Tragen:] sind mit allen Stollenwerk-Fahrgestellen kompatibel!
\end{inparadesc}   
    \item[Max. Belastbarkeit:] ---\Kg*
    \item[Homepage:] \url{http://www.stollenwerk-koeln.de}
%     \item[Hygienische Wiederaufbereitung:]
\end{compactdesc}
% {\TakeNotesLarge}
}{}{}{}{}

\Layout{37}{Abbildung}
{}
{}
{./images/pallinger-christoph-aass/Trage_32970_v2-AASS-0176mm.jpg}
{Krankentrage der Firma Stollenwerk mit Fahrgestell.}
{Ch. Pallinger}

\Layout{2}{Funktionen}
{%
~

{\TakeNotesLarge}
}{}{}{}{}

\Layout{2}{Anwendung}
{%
~

{\TakeNotesLarge}
}{}{}{}{}

\Layout{2}{Hygienische Wiederaufbereitung}
{%
~

{\TakeNotesLarge}
}{}{}{}{}

\Layout{2}{\TxGefahren}
{%
~

{\TakeNotesLarge}
}{}{}{}{}




\subsubsection{Krankentrage mit Fahrgestell der Fa. Ferno}
% TODO : Neufassung
\TODO{GABS}{yyyy-mm-dd}{Krankentrage Ferno: Fehlender
Text.}{}
\Layout{2}{\TxBeschreibung}
{%
~

{\TakeNotesLarge}
}{}{}{}{}

\Layout{2}{Funktionen}
{%
~

{\TakeNotesLarge}
}{}{}{}{}

\Layout{2}{Anwendung}
{%
~

{\TakeNotesLarge}
}{}{}{}{}

\Layout{2}{Hygienische Wiederaufbereitung}
{%
~

{\TakeNotesLarge}
}{}{}{}{}

\Layout{2}{\TxGefahren}
{%
~

{\TakeNotesLarge}
}{}{}{}{}




\section{Geräte zum Retten von Patienten}

\subsection{Schaufeltrage}


\subsection{Spineboard}


\subsection{Rettungskorsett}


   
\section{Schienungsmaterial}
\label{topic:schienungsmaterial}





%%%
\nocite{RueckerBildatlas1}